% RESEARCH_DRAFT_SHA256: 2a95c59293ef2017e0df4980df3b96fa3b7c849375dc7a603a391fa9db16e702
% Layout and source-reference wrappers expanded for the web renderer.
\documentclass[11pt]{article}
\usepackage[margin=25mm]{geometry}
\usepackage{amsthm}
\usepackage{fontspec,unicode-math}
\setmainfont{Latin Modern Roman}
\setmathfont{Latin Modern Math}
\usepackage{amsmath,amssymb,mathtools}
\usepackage{xurl,hyperref}
\hypersetup{colorlinks=true,urlcolor=blue,linkcolor=blue}
\setcounter{secnumdepth}{0}
\setlength{\parindent}{0pt}
\setlength{\parskip}{5pt}
\setlength{\emergencystretch}{3em}
\newcommand{\sref}[2]{\hyperlink{source:#1:#2}{[S#2]}}
\newcommand{\cataloguescope}[1]{\paragraph{Source question.}#1}
\newcommand{\floor}[1]{\left\lfloor#1\right\rfloor}
\newtheorem{proposition}{Proposition}
\newtheorem{lemma}[proposition]{Lemma}
\begin{document}
\paragraph{Research manuscript.} OpenAI \texttt{gpt-6-astra}, reasoning \texttt{ultra}; 11 October 2026.
\paragraph{Preserved source body.} The complete catalogue mathematical text and sources follow with layout and source links expanded, including source qualifications and its historical preparation notes. New research begins at the section headed Research attempt.

\section{Recognition and Characterization of Leaf Powers}\label{statement}

\subsection{Definitions and mathematical statement}

\paragraph{Shared notation.}Standard number systems are denoted by $\mathbb N,\mathbb Z,\mathbb Q,\mathbb R,\mathbb C$; index ranges, computational models and specialized notation follow the cited problem source and the definitions below.


A graph is a $k$-leaf power if there is a tree whose leaves are its vertices and two vertices are adjacent exactly when their distance in the tree is at most $k$. Such a tree is a $k$-leaf root.
\paragraph{Statement recorded in the supplied source.}
For each $k\geq5$, characterize $k$-leaf powers and determine the complexity of recognizing them. \href{https://doi.org/10.4171/owr/2006/07}{[S2]}
\paragraph{Source status.}
The supplied catalogue labels this record partially solved. Its source-status metadata was checked against the supplied snapshot on 2026-10-11. This import does not independently establish current open status or mathematical correctness. The catalogue literature assessment is dated 2026-08-21.
\subsection{Short English statement}
Characterize and recognize leaf-power graphs for every distance threshold at least five.
\subsection{Sources}
\begin{itemize}
\item[S1] Ulam AI and the UnsolvedMath Contributors, UnsolvedMath: A Curated Collection of Open Mathematics Problems, catalogue entry. Catalogue curation and metadata: CC BY 4.0. The supplied source question is retained; definitions and the English summary are editorial additions.
\url{https://huggingface.co/datasets/ulamai/UnsolvedMath}.
\item[S2] Algorithmic Graph Theory (2006), Oberwolfach Reports, DOI 10.4171/owr/2006/07, pp. 379--460
\url{https://doi.org/10.4171/owr/2006/07}.
\item[S3] Lafond recognition theorem
\url{https://arxiv.org/abs/2110.15421}.
\item[S4] Dupré la Tour et al. obstruction theorem
\url{https://arxiv.org/abs/2407.02412}.
\textit{Underlying Oberwolfach Reports material retains its source-specific licence, including CC BY-SA 4.0 where applicable; see THIRD\_PARTY\_NOTICES.md.}
\end{itemize}
\paragraph{Preparation and review.}
Original question attributed to the cited source and UnsolvedMath contributors. Codex (GPT-6; exact serving revision not exposed) assisted with definitions, English summaries, source inspection and formatting. No new human mathematical authorship or independent proof verification is asserted. Catalogue identity checked against the complete supplied registry and existing repository statements; known flagged repetitions and obvious identical conjectures excluded. Full mathematical-equivalence/literature review has not been performed. Statement preparation and final review followed contribution prompts C and D; automated checks verify structure and identity, not open status.
\label{end:statement}

\section{Research attempt}
\subsection{Literature, source context and scope}
The full original catalogue body and source qualifications are preserved above. The relevant primary workshop passages were read in the cached publisher report. Historical claims and subsequent attributed results are distinguished from the mathematical work below. A restricted example, standard lemma or known theorem is not counted as new progress on the entire original target. Missing assumptions and repeated or rhetorical extracts are not silently turned into different solved problems.

\subsection{Compressing a leaf-root certificate}
The historical report asks about recognition for fixed $k\ge5$. Subsequent primary literature states a polynomial-time algorithm for each fixed $k$, with a very large dependence on $k$, and separately rules out a finite forbidden-induced-subgraph characterization even within strongly chordal graphs for $k\ge5$. These are attributed literature results; neither proof is reproduced here. We investigate the simpler certificate-enumeration route and its precise limitation.
\begin{proposition}
For a graph on $n\ge2$ vertices and an integer $k\ge1$, existence of a unit-edge $k$-leaf root implies existence of a weighted-tree certificate with the same labeled leaves, at most $n-2$ other vertices, at most $2n-3$ edges, and integer edge lengths in $\{1,\ldots,k+1\}$. The converse also holds after subdividing weighted edges. The certificate uses $O(n\log n+n\log(k+1))$ bits and can be verified in polynomial time in this size.
\end{proposition}
\begin{proof}
In a unit-edge root, suppress every degree-two nonleaf vertex, replacing its incident path by one edge whose length is that path's positive integer length. Leaf distances are unchanged. Every remaining nonleaf has degree at least three. If there are $b$ nonleaves, the degree sum gives $n+3b\le2(n+b-1)$, hence $b\le n-2$ and the number of edges $n+b-1$ is at most $2n-3$.
Replace each edge length $w$ by $\min(w,k+1)$. A leaf-to-leaf path containing a length greater than $k$ remains longer than $k$ afterward, because that edge alone still has length $k+1$. Every path containing only lengths at most $k$ is unchanged. Thus the distance-at-most-$k$ relation is preserved. Conversely subdividing each weighted edge into a path of its integer length realizes the same distances with unit edges and introduces no extra leaves.
Encode the topology using $O(n)$ edges with vertex labels of $O(\log n)$ bits, and encode the $O(n)$ bounded lengths with $O(\log(k+1))$ bits each. Pairwise leaf distances can be summed in binary along the unique paths and compared with $k$; they have at most $O(\log n+\log(k+1))$ bits. Checking every pair and the tree structure is therefore polynomial in the certificate size, without expanding long edges.
\end{proof}
For $n=1$, the sole graph is handled separately using the conventional one-vertex tree as a trivial leaf root; this boundary case is not needed by the suppression count. The problem's actual $k\ge5$ lies inside the proposition's positive-integer range.
\subsection{Verification and remaining gaps}
For a weighted star with arm lengths $1,2,6,7$ and $k=5$, capping at 6 preserves exactly the same leaf adjacencies: only the first two leaves are adjacent. All six pair distances were checked. Enumeration of the compressed trees still permits exponentially many topologies and length assignments, so the step from a short verifiable witness to a polynomial recognition algorithm is invalid. The later fixed-$k$ algorithm is not derived by this argument. No new recognition bound, uniform parameter dependence or forbidden-structure theorem is established.

\section{References and source attribution}
Primary workshop report: \url{https://doi.org/10.4171/owr/2006/07}. Relevant passages read: 2006/7 p. 451, leaf powers, leaf roots and the historical recognition/characterization question.

\url{https://arxiv.org/html/2110.15421v1}. Read primary Lafond paper statement of polynomial recognition for each fixed k and its large parameter dependence; no independent proof verification.

\url{https://arxiv.org/abs/2407.02412}. Read primary abstract stating absence of a finite forbidden-induced-subgraph characterization within strongly chordal graphs for every k>=5; proof not independently checked.

All retrievals were checked on 11 October 2026. Original question attribution remains in the preserved Sources section. External theorems are attributed and not claimed as new discoveries or fully verified merely by citation.


\paragraph{Integration review.} Codex (GPT-6; exact serving revision not exposed) reviewed the written arguments and expanded layout and source-reference wrappers. Repository metadata and the production web renderer were checked using Python 3.13 and JavaScript with MathJax 3.2.2. This model-assisted review is not independent human proof verification; no PDF compilation or formal proof checking is asserted.

\section{Completion Estimate}
\noindent\textbf{COMPLETION: 0\%}
No substantiated new progress toward the entire original target is claimed. The worked results and failed route are delimited above; a known result, illustrative model, or auxiliary obstruction is not counted as original-target completion.

\section{Final status and attribution}
\textbf{Attempt status: unresolved.} This describes this research attempt, not a new ruling on the source catalogue's status or on external literature claims.
The mathematics and this self-audit were generated by OpenAI \texttt{gpt-6-astra}, reasoning \texttt{ultra}, on 11 October 2026. Codex (GPT-6; exact serving revision not exposed) orchestrated the batch and contributed formatting/repository review. No human mathematical authorship, independent human proof checking, or formal verification is claimed. The historical preparation paragraphs in the reproduced catalogue body concern that body only.

Contribution prompt A was applied to the complete definition and supplied source record, followed by the adversarial checks of prompt B. Local reproduction files are \texttt{batch016/checks.py}, \texttt{computations.json}, and the per-ID audit JSON. They were not public links at preparation time. Numerical checks are finite and do not prove unrestricted results. No LaTeX engine was found locally; compilation was not performed by this research agent.

\end{document}
